%! Exponential Function Manipulation — Unit 4, Lesson 4.4 — Guided Notes (Answer Key)
\documentclass[10pt]{article}
\usepackage{precalculus-article}
\usepackage{precalculus-key}   % pulls in -boxes; do NOT also load -boxes
\newcommand{\TallMath}[1]{$\displaystyle #1\rule[-1.4em]{0pt}{3.2em}$}

\begin{document}

\pageheader{Unit 4, Lesson 4.4}{Guided Notes --- Answer Key}
\namedateperiod

\vspace{0.08in}

\begin{objectivebox}
By the end of this lesson, I will be able to\ldots
\begin{enumerate}[itemsep=2pt]
  \item Apply the \ans{product} property to rewrite a horizontal translation $b^{x+k}$ as a
        vertical dilation $a \cdot b^x$. \hfill\textit{(EK 2.4.A.1)}
  \item Apply the \ans{power} property to rewrite a horizontal dilation $b^{cx}$ as a
        \ans{base} change. \hfill\textit{(EK 2.4.A.2)}
  \item Simplify expressions with \ans{negative} exponents and \ans{fractional} exponents
        (roots). \hfill\textit{(EK 2.4.A.3, A.4)}
\end{enumerate}
\end{objectivebox}

\vspace{0.06in}

\begin{vocabbox}
\termblanklong{Product property of exponents}
\termblanklong{Power property of exponents}
\termblanklong{Negative exponent}
\termblanklong{Fractional exponent}
\termblanklong{$k$th root}
\termblanklong{Equivalent form}
\termblanklong{Vertical dilation}
\termblanklong{Horizontal translation}
\end{vocabbox}

\vspace{0.06in}

\begin{hookbox}
Two functions are on the board:
\[
  f(x) = 2^{x+3} \qquad\qquad g(x) = 8 \cdot 2^x
\]
Evaluate each at $x = 0$: $f(0) = $ \ans{8}, $g(0) = $ \ans{8}. \\[4pt]
Evaluate each at $x = 1$: $f(1) = $ \ans{16}, $g(1) = $ \ans{16}. \\[4pt]
These functions \ans{are} (are / are not) equivalent because $2^{x+3} = 2^x \cdot 2^{{\color{keyred}\mathbf{3}}} = $ \ans{8} $\cdot\, 2^x$.
\end{hookbox}

\vspace{0.06in}

\begin{notesbox}{LO 2.4.A --- Product Property (EK 2.4.A.1)}
\textbf{Product Property:} \quad $b^m \cdot b^n = $ \ans{$b^{m+n}$}

\medskip
\textbf{Graphical connection:} A horizontal translation $f(x) = b^{x+k}$ is equivalent to a
vertical dilation $f(x) = $ \ans{$b^k$} $\cdot\, b^x$, where $a = $ \ans{$b^k$}.

\medskip
\textbf{Examples --- rewrite in the form $ab^x$:}

\renewcommand{\arraystretch}{2.2}
\begin{tabular}{@{} l @{\quad$=$\quad} l @{\quad$=$\quad} l @{}}
$5^{x+2}$ & $5^x \cdot 5^{{\color{keyred}\mathbf{2}}}$ & \ans{25} $\cdot\, 5^x$ \\
$3^{x-1}$ & $3^x \cdot 3^{{\color{keyred}\mathbf{-1}}}$ & \ans{$\frac{1}{3}$} $\cdot\, 3^x$ \\
$4 \cdot 2^x$ & $2^{{\color{keyred}\mathbf{2}}} \cdot 2^x$ & $2^{x + {\color{keyred}\mathbf{2}}}$ \\
\end{tabular}
\end{notesbox}

\vspace{0.06in}

\begin{notesbox}{LO 2.4.A --- Power Property (EK 2.4.A.2)}
\textbf{Power Property:} \quad $(b^m)^n = $ \ans{$b^{mn}$}

\medskip
\textbf{Graphical connection:} A horizontal dilation $f(x) = b^{cx}$ is equivalent to a base
change $f(x) = ($ \ans{$b^c$} $)^x$, where the new base is $b^c$.

\medskip
\textbf{Examples --- rewrite as a single base to the $x$ power:}

\renewcommand{\arraystretch}{2.2}
\begin{tabular}{@{} l @{\quad$=$\quad} l @{\quad$=$\quad} l @{}}
$2^{3x}$   & $(2^{{\color{keyred}\mathbf{3}}})^x$ & \ans{8}$^x$ \\
$9^{x/2}$  & $(9^{{\color{keyred}\mathbf{1/2}}})^x$ & \ans{3}$^x$ \\
$4^{x/2}$  & $(4^{{\color{keyred}\mathbf{1/2}}})^x$ & \ans{2}$^x$ \\
\end{tabular}
\end{notesbox}

\vspace{0.06in}

\begin{notesbox}{LO 2.4.A --- Negative \& Fractional Exponents (EK 2.4.A.3, A.4)}
\renewcommand{\arraystretch}{1.5}
\begin{tabular}{@{} >{\bfseries}p{0.38\linewidth} l @{}}
Negative exponent rule:     & $b^{-n} = \dfrac{1}{b^n}$ \ans{(reciprocal)} \\
Fractional exponent rule:   & $b^{1/k} = $ \ans{$\sqrt[k]{b}$} \\
General fractional exponent: & $b^{m/n} = (\sqrt[n]{b})^m$, exponent $= $ \ans{$m$} \\
\end{tabular}

\medskip
\textbf{Decay connection:} $f(x) = 2^{-x} = \left(\dfrac{1}{2}\right)^x$ \ans{(base flipped)}

\medskip
\textbf{Evaluate:}
\quad $2^{-3} = $ \ans{$\dfrac{1}{8}$}
\quad $8^{1/3} = $ \ans{2}
\quad $4^{3/2} = $ \ans{8}
\end{notesbox}

\end{document}
